\begin{textsample}{NEG Dim 3 – global\_north\_2025\_04 – Score 1.00 – global\_north\_2025\_04/124112364.txt}  \label{ex:f3_neg_416}
The Palestine Solidarity Network Aotearoa is demanding the New Zealand government justify its absence from submitters to the International Court of Justice hearings at the Hague into Israel blocking vital supplies entering Gaza.
The ICJ ' s ongoing investigation into Israeli genocide in Gaza is now considering the illegality of Israel cutting off all food, water, fuel, medicine and other essential aid entering Gaza since early March.
Countries submitting include the UK, Spain, Belgium and Malaysia.
New Zealand is not on the list for making a submission.
PSNA Co-Chair Maher Nazzal says the New Zealand government has gone completely silent on Israeli atrocities in Gaza."A year ago, the Prime Minister and Foreign Minister were making statements about how Israel must comply with international law.""They carefully avoided blaming Israel for doing anything wrong, but they issued strong warnings, such as telling Israel that it should not attack the city of Rafah.""Israel then bombed Rafah flat.
The New Zealand response was to @ @ @ @ @ @ @ @ @ @ open about driving Palestinians out of Gaza, so Israel can build Israeli settlements there.
And they are just as open about using starvation as a weapon.""Our government says and does nothing.
Chris Luxon had nothing to say about Gaza when he met British Prime Minister Keir Stamer in London earlier in the month.
Yet Israel is perpetuating the holocaust of the 21st Century under the noses of both Prime Ministers."Advertisement - scroll to continue reading Maher Nazzal says that it is deeply disappointing that a nation which so proudly invokes its history of standing against apartheid and of championing nuclear disarmament, chooses to not even appear on the sidelines of the ICJ ' s legal considerations."New Zealand can not claim to stand for a rules-based international order while selectively avoiding the rules when it comes to Palestine.""We want the New Zealand government to urgently explain to the public its absence from the ICJ hearings.
We need it to commit to participating in all future international legal @ @ @ @ @ @ @ @ @ @ obligations to impose sanctions on Israel to force its withdrawal from the Palestinian Occupied Territory.""If even small countries, such as Vanuatu, can commit their meagre resources to go to make a case to the ICJ, then surely our government can at the very least do the same. ' See here for the official list of countries and other organisations submitting to the ICJ Using Scoop for work?
Scoop is free for personal use, but you ' ll need a licence for work use.
This is part of our Ethical Paywall and how we fund Scoop.
Join today with plans starting from less than \$3 per week, plus gain access to exclusive Pro features.

% matched lemmas: 
\end{textsample}
